\documentclass[11pt]{article}

\usepackage[margin=1in]{geometry}
\usepackage{amsmath}
\usepackage{amssymb}
\usepackage{array}
\usepackage{booktabs}
\usepackage[hidelinks]{hyperref}
\usepackage{microtype}

\newcommand{\Uone}{\mathrm{U}(1)}
\newcommand{\SUN}[1]{\mathrm{SU}(#1)}
\newcommand{\ReTr}{\operatorname{Re}\operatorname{Tr}}
\newcommand{\ee}{\mathrm{e}}
\newcommand{\Vlat}{\mathcal{V}}
\newcommand{\Nmeas}{N_{\mathrm{meas}}}

\newenvironment{interpretation}
  {\begin{quote}\small\noindent\textbf{Interpretation.}\ }
  {\end{quote}}

\title{Lattice Monte Carlo Foundations\\
\large The Two-Dimensional Ising Model and Pure-Gauge Calculations}
\author{QCD Prediction Project}
\date{}

\begin{document}

\maketitle

\begin{abstract}
This document explains the mathematical operations in two related lattice
Monte Carlo tracks.  The two-dimensional Ising model supplies a small,
exactly enumerable test of local Metropolis updates, thermalization, and
correlated measurements.  The pure-gauge track then applies the same numerical
ideas to compact \(\Uone\), pure \(\SUN{2}\), and pure \(\SUN{3}\) models with
periodic boundaries and the Wilson gauge action.  Each expression is connected
to the data it consumes, the quantity it produces, and its algorithmic role.
Dynamical fermions are outside this calculation.
\end{abstract}

\tableofcontents

\section{Calculation at a glance}

Both tracks begin with the same validated lattice geometry, deterministic
random streams, sampling schedule, and statistical analysis.  The Ising track
stores spins on sites and is the first end-to-end validation problem.  The
gauge track stores group elements on links and has the following dependency
chain:

\begin{enumerate}
  \item Convert lattice extents into site indices and periodic neighbors.
  \item Allocate and initialize one gauge link per site and direction.
  \item Form plaquettes and evaluate the Wilson gauge action.
  \item Update links with a Markov-chain Monte Carlo algorithm.
  \item Discard thermalization, then measure plaquettes and Wilson loops.
  \item Estimate autocorrelation-aware statistical uncertainties.
  \item Convert Wilson-loop ratios into an effective static potential.
  \item Validate structural identities and compare independent chains.
\end{enumerate}

The output of one stage is the input of the next:

\begin{center}
\begin{tabular}{@{}p{0.18\textwidth}p{0.31\textwidth}p{0.38\textwidth}@{}}
\toprule
\textbf{Stage} & \textbf{Consumes} & \textbf{Produces} \\
\midrule
Geometry & Extents and boundaries & Indices, strides, neighbors \\
Gauge field & Geometry and initial-state policy & Links \(U_\mu(x)\) \\
Action & Links and coupling \(\beta\) & Plaquettes and \(S_W[U]\) \\
Update & Current field and random numbers & Next Markov-chain state \\
Measurement & Thermalized field samples & \(\bar P\) and \(W(r,t)\) \\
Statistics & Correlated measurement history & Means, covariance, errors \\
Potential & Correlated Wilson loops & \(aV_{\mathrm{eff}}(r,t)\), then \(V(r)\) \\
\bottomrule
\end{tabular}
\end{center}

\section{Notation}

\begin{center}
\begin{tabular}{@{}>{$}l<{$}p{0.72\textwidth}@{}}
\toprule
\multicolumn{1}{c}{\textbf{Symbol}} & \textbf{Meaning} \\
\midrule
D & Number of Euclidean lattice dimensions. \\
x=(x_0,\ldots,x_{D-1}) & Integer coordinate of one lattice site. \\
L_\mu & Number of sites in direction \(\mu\). \\
\Vlat & Total number of lattice sites. A calligraphic symbol distinguishes
volume from the potential \(V(r)\). \\
\mu,\nu & Direction indices in \(\{0,\ldots,D-1\}\). \\
\hat\mu & Unit displacement in direction \(\mu\). \\
U_\mu(x) & Gauge link from \(x\) to \(x+\hat\mu\). \\
G & Gauge group: initially \(\Uone\), then \(\SUN{2}\) or \(\SUN{3}\). \\
N & Matrix dimension for \(\SUN{N}\); QCD uses \(N=3\). \\
s_x & Ising spin at site \(x\), with \(s_x\in\{-1,+1\}\). \\
J & Ising nearest-neighbor interaction strength. \\
h & Ising external magnetic field; the first implementation uses \(h=0\). \\
\beta_{\mathrm{th}} & Thermodynamic inverse temperature \(1/(k_{\mathrm B}T)\). \\
\beta & Bare gauge coupling parameter appearing inside the Wilson action. \\
a & Lattice spacing; \(aV\) is a dimensionless potential. \\
r,t & Spatial and Euclidean-time extents in lattice units. \\
\Nmeas & Number of saved measurements in a chain. \\
\bottomrule
\end{tabular}
\end{center}

\section{Step 1: construct the lattice}

\subsection{Allowed sites and volume}

For periodic extents \(L_0,\ldots,L_{D-1}\), the set of sites is
\begin{equation}
  \Lambda =
  \left\{x=(x_0,\ldots,x_{D-1})\;\middle|\;
  0\leq x_\mu<L_\mu\right\}.
  \label{eq:lattice-sites}
\end{equation}

\begin{interpretation}
Each component \(x_\mu\) is an integer. Equation~\eqref{eq:lattice-sites}
defines which coordinates are legal before any fields are allocated.
\end{interpretation}

The total number of sites is
\begin{equation}
  \Vlat = |\Lambda| = \prod_{\mu=0}^{D-1}L_\mu.
  \label{eq:lattice-volume}
\end{equation}

\begin{interpretation}
The program must check that every \(L_\mu>0\) and that the product in
Equation~\eqref{eq:lattice-volume} does not overflow the index type.
\(\Vlat\) determines the site allocation size.
\end{interpretation}

\subsection{Flattened storage}

For row-major storage, define strides and the flattened site index by
\begin{equation}
  s_0=1,\qquad
  s_\mu=\prod_{\nu=0}^{\mu-1}L_\nu,\qquad
  i(x)=\sum_{\mu=0}^{D-1}x_\mu s_\mu.
  \label{eq:site-index}
\end{equation}

\begin{interpretation}
\(s_\mu\) is the number of contiguous entries skipped when \(x_\mu\)
increases by one. The result \(i(x)\) lies in
\(0\leq i(x)<\Vlat\) and is used to address site-owned data.
\end{interpretation}

Because every site owns one outgoing link in each direction, a convenient link
index is
\begin{equation}
  \ell(x,\mu)=D\,i(x)+\mu,
  \qquad 0\leq\ell<D\Vlat.
  \label{eq:link-index}
\end{equation}

\begin{interpretation}
Equation~\eqref{eq:link-index} maps a site and direction to one flat gauge-field
allocation containing \(D\Vlat\) group elements.
\end{interpretation}

\subsection{Periodic neighbors}

Forward and backward neighbors change only coordinate \(\mu\):
\begin{equation}
\begin{aligned}
  (x+\hat\mu)_\nu &=
  \begin{cases}
    (x_\mu+1)\bmod L_\mu, & \nu=\mu,\\
    x_\nu,                & \nu\neq\mu,
  \end{cases}
  \\[1ex]
  (x-\hat\mu)_\nu &=
  \begin{cases}
    (x_\mu+L_\mu-1)\bmod L_\mu, & \nu=\mu,\\
    x_\nu,                      & \nu\neq\mu.
  \end{cases}
  \label{eq:periodic-neighbors}
\end{aligned}
\end{equation}

\begin{interpretation}
The added \(L_\mu\) in the backward expression keeps the dividend nonnegative
before the modulo operation. These neighbors are reused in plaquettes, staples,
Wilson paths, and boundary tests.
\end{interpretation}

\paragraph{Step output.}
A validated lattice shape, a coordinate/index round trip, and forward/backward
neighbor operations.

\section{Toy-model track: the two-dimensional Ising model}

The Ising model is not a gauge theory and is not itself QCD.  It is a compact
test problem for the same Markov-chain, random-number, measurement, and error
analysis machinery needed later.  The first implementation uses a periodic
\(L_x\times L_y\) square lattice and one spin
\begin{equation}
  s_x\in\{-1,+1\}
  \label{eq:ising-spin}
\end{equation}
at every site.

\subsection{Hamiltonian and bond convention}

Count each positive-axis bond once and define
\begin{equation}
  E[s]
  =-J\sum_{x\in\Lambda}
      \left[s_xs_{x+\hat 0}+s_xs_{x+\hat 1}\right]
    -h\sum_{x\in\Lambda}s_x.
  \label{eq:ising-energy}
\end{equation}
For positive \(J\), aligned neighboring spins lower the energy.  Equation
\eqref{eq:ising-energy} contains \(2\Vlat\) bond terms, and periodic neighbors
are evaluated with Equation~\eqref{eq:periodic-neighbors}.

On a \(2\times2\) periodic lattice, the forward and backward neighbor in one
axis have the same coordinate.  They must nevertheless appear twice in the
local-neighbor multiset because they represent two periodic bond occurrences.
Removing the duplicate changes the model and invalidates the exact reference
values below.

The equilibrium probability of a configuration is
\begin{equation}
  \pi[s]=\frac{1}{Z}\exp\!\left[-\beta_{\mathrm{th}}E[s]\right],
  \qquad
  Z=\sum_{\{s\}}\exp\!\left[-\beta_{\mathrm{th}}E[s]\right].
  \label{eq:ising-boltzmann}
\end{equation}
It is often convenient to use the dimensionless coupling
\(K=\beta_{\mathrm{th}}J\).  Here \(\beta_{\mathrm{th}}=1/(k_{\mathrm B}T)\)
is an inverse temperature.  It is conceptually different from the bare gauge
parameter \(\beta\) in the Wilson action.

\subsection{Local spin-flip energy}

Propose changing only the spin at site \(x\):
\begin{equation}
  s_x'=-s_x.
  \label{eq:ising-flip}
\end{equation}
Every term not containing \(s_x\) cancels in
\(\Delta E=E[s']-E[s]\).  If \(\mathcal N(x)\) is the four-entry multiset
containing the forward and backward neighbor in both axes, the affected energy
before the flip is
\begin{equation}
  E_{\mathrm{old},x}
  =-J s_x\sum_{y\in\mathcal N(x)}s_y-hs_x,
  \label{eq:ising-old-local}
\end{equation}
and after the flip it is
\begin{equation}
  E_{\mathrm{new},x}
  =+J s_x\sum_{y\in\mathcal N(x)}s_y+hs_x.
  \label{eq:ising-new-local}
\end{equation}
Subtracting gives the local formula
\begin{equation}
  \boxed{\Delta E
  =2s_x\left(J\sum_{y\in\mathcal N(x)}s_y+h\right)}.
  \label{eq:ising-delta-energy}
\end{equation}
This formula should be tested against a deliberately slow calculation of the
full energy before and after every possible spin flip on small lattices.

\subsection{Metropolis update and sweep definition}

For a symmetric flip proposal, accept with probability
\begin{equation}
  P_{\mathrm{accept}}
  =\min\left(1,\exp[-\beta_{\mathrm{th}}\Delta E]\right).
  \label{eq:ising-acceptance}
\end{equation}
For \(u\sim\operatorname{Uniform}[0,1)\), the numerically stable decision is
\begin{equation}
  \log u<-\beta_{\mathrm{th}}\Delta E.
  \label{eq:ising-log-acceptance}
\end{equation}
When \(\Delta E\leq0\), accept without evaluating the logarithm.  One sweep is
exactly \(\Vlat\) attempted spin flips.  The program must record whether those
attempts use a fixed site order, a new random permutation each sweep, or random
sites sampled with replacement; these policies are not identical finite-chain
procedures even though each can target the same equilibrium distribution.

A minimal run is
\begin{verbatim}
initialize all spins and a deterministic random stream

repeat N_warmup times:
    attempt one sweep of local spin flips

repeat N_measurements times:
    repeat N_spacing times:
        attempt one sweep of local spin flips
    record energy, magnetization, acceptance count, and sweep number
\end{verbatim}
Warm-up length is not proven adequate merely because it was configured.
Compare ordered and random initial states, inspect histories, and repeat with a
longer warm-up and measurement spacing.

\subsection{Measurements}

For \(\Vlat=L_xL_y\), define total and per-site energy by
\begin{equation}
  E=E[s],
  \qquad
  e=\frac{E}{\Vlat},
  \label{eq:ising-energy-density}
\end{equation}
and total and per-site magnetization by
\begin{equation}
  M=\sum_{x\in\Lambda}s_x,
  \qquad
  m=\frac{M}{\Vlat}.
  \label{eq:ising-magnetization}
\end{equation}
At zero field and finite volume, spin-flip symmetry gives
\(\langle m\rangle=0\) when both magnetized sectors are sampled.  Therefore
record \(m\), \(|m|\), and \(m^2\), rather than interpreting a near-zero mean
magnetization as an absence of ordering.  Also record the acceptance fraction
\begin{equation}
  f_{\mathrm{accept}}
  =\frac{N_{\mathrm{accepted}}}{N_{\mathrm{attempted}}}.
  \label{eq:ising-acceptance-fraction}
\end{equation}
Measurements from consecutive sweeps are correlated and must use the
autocorrelation or blocking treatment in Step~7.

\subsection{Exact \texorpdfstring{\(2\times2\)}{2x2} enumeration}

For \(h=0\), enumerate all \(2^{\Vlat}=16\) configurations of a
\(2\times2\) lattice.  With the bond convention in Equation
\eqref{eq:ising-energy}, the energy levels and degeneracies are
\begin{center}
\begin{tabular}{@{}ccc@{}}
\toprule
Energy & Degeneracy & Configuration class \\
\midrule
\(-8J\) & 2 & all spins aligned \\
\(0\) & 12 & one-versus-three or adjacent two-versus-two \\
\(+8J\) & 2 & checkerboard \\
\bottomrule
\end{tabular}
\end{center}
Consequently,
\begin{equation}
  Z_{2\times2}
  =2\ee^{8K}+12+2\ee^{-8K}
  =4\left[\cosh(8K)+3\right],
  \label{eq:ising-2x2-partition}
\end{equation}
and the exact per-site energy is
\begin{equation}
  \langle e\rangle_{2\times2}
  =-2J\frac{\sinh(8K)}{\cosh(8K)+3}.
  \label{eq:ising-2x2-energy}
\end{equation}
The exact absolute magnetization per site is
\begin{equation}
  \langle |m|\rangle_{2\times2}
  =\frac{\ee^{8K}+2}{2[\cosh(8K)+3]},
  \qquad
  \langle m\rangle_{2\times2}=0.
  \label{eq:ising-2x2-magnetization}
\end{equation}
At \(K=0\), these expressions give \(\langle e\rangle=0\) and
\(\langle|m|\rangle=3/8\).  As \(K\to\infty\), they approach
\(\langle e\rangle=-2J\) and \(\langle|m|\rangle=1\).

The enumerator should construct configurations independently of the Monte
Carlo code, evaluate Equation~\eqref{eq:ising-energy} directly, and accumulate
weighted observables using Equation~\eqref{eq:ising-boltzmann}.  It is a test
oracle, not a production sampling algorithm.

\subsection{Ising validation gates}

Before using the toy simulation as evidence that the sampling infrastructure
works, require all of the following:
\begin{enumerate}
  \item every local \(\Delta E\) equals a full-energy recomputation;
  \item the enumerator reproduces the degeneracies and exact formulas above;
  \item the same seed and update policy reproduce the same CPU history;
  \item ordered and random starts agree after adequate warm-up;
  \item Monte Carlo estimates agree with exact \(2\times2\) values within
        autocorrelation-aware uncertainty;
  \item uncertainty grows when fewer effectively independent samples are used;
        and
  \item measurements retain chain identity and sweep number.
\end{enumerate}

\subsection{Relationship to the gauge track}

\begin{center}
\begin{tabular}{@{}p{0.21\textwidth}p{0.31\textwidth}p{0.36\textwidth}@{}}
\toprule
Concept & Ising track & Pure-gauge track \\
\midrule
Degrees of freedom & Spins \(s_x\) on sites & Group elements \(U_\mu(x)\) on links \\
Statistical weight & \(\exp[-\beta_{\mathrm{th}}E]\) & \(\exp[-S_W]\), with \(\beta\) inside \(S_W\) \\
Local change & \(\Delta E\) from neighboring spins & \(\Delta S\) from neighboring plaquettes \\
First observables & \(e\), \(m\), \(|m|\) & Plaquette and Wilson loops \\
Shared machinery & Geometry, RNG, sweeps, histories, autocorrelation, uncertainty & Same \\
\bottomrule
\end{tabular}
\end{center}

\section{Step 2: initialize the gauge field}

A directed link stores an element of the chosen gauge group:
\begin{equation}
  U_\mu(x)\in G,\qquad
  G\in\{\Uone,\SUN{2},\SUN{3}\}.
  \label{eq:gauge-link}
\end{equation}

For an \(\SUN{N}\) link, group membership requires
\begin{equation}
  U_\mu^\dagger(x)U_\mu(x)=I,
  \qquad
  \det U_\mu(x)=1.
  \label{eq:su-constraints}
\end{equation}

\begin{interpretation}
\(U_\mu(x)\) transports a color vector from \(x\) to \(x+\hat\mu\).
Traversing the link backward uses \(U_\mu^\dagger(x)\). The identities in
Equation~\eqref{eq:su-constraints} are also inexpensive numerical diagnostics.
\end{interpretation}

A cold start uses \(U_\mu(x)=I\). A hot start draws a random valid group
element for each link. Both starts should approach the same equilibrium
distribution after sufficient thermalization.

\paragraph{Step output.}
A gauge field containing \(D\Vlat\) valid link values and reproducibly seeded
random-number streams.

\section{Step 3: form plaquettes and the Wilson action}

\subsection{Oriented plaquette}

The elementary loop in the positive \(\mu\)-\(\nu\) orientation is
\begin{equation}
  U_{\mu\nu}(x)=
  U_\mu(x)\,
  U_\nu(x+\hat\mu)\,
  U_\mu^\dagger(x+\hat\nu)\,
  U_\nu^\dagger(x).
  \label{eq:plaquette}
\end{equation}

\begin{interpretation}
Read Equation~\eqref{eq:plaquette} from left to right as four path segments:
forward in \(\mu\), forward in \(\nu\), backward in \(\mu\), and backward in
\(\nu\). The path returns to \(x\). Its traced value is gauge invariant even
though an individual link is not.
\end{interpretation}

\subsection{Wilson gauge action}

For pure \(\SUN{N}\), use
\begin{equation}
  S_W[U]=
  \beta\sum_x\sum_{\mu<\nu}
  \left(
    1-\frac{1}{N}\ReTr U_{\mu\nu}(x)
  \right).
  \label{eq:wilson-action}
\end{equation}

\begin{interpretation}
The inner sum counts every unoriented plaquette once. A plaquette close to the
identity has a small contribution; rough field configurations have a larger
action. The convention in Equation~\eqref{eq:wilson-action} must be used
consistently by the global action, local update, and observables.
\end{interpretation}

The corresponding ensemble expectation of an observable \(\mathcal O\) is
\begin{equation}
  \langle\mathcal O\rangle =
  \frac{1}{Z}\int\mathcal D U\;
  \mathcal O[U]\ee^{-S_W[U]},
  \qquad
  Z=\int\mathcal D U\;\ee^{-S_W[U]}.
  \label{eq:expectation}
\end{equation}

\begin{interpretation}
\(\mathcal D U\) means integration over every gauge link using the invariant
group measure. \(Z\) normalizes the probability distribution. The program does
not enumerate this integral; it estimates it with correlated samples from a
Markov chain whose target weight is proportional to \(\ee^{-S_W[U]}\).
\end{interpretation}

\paragraph{Step output.}
A global action definition, a local plaquette primitive, and a target
probability distribution for the sampler.

\section{Step 4: update one gauge link}

\subsection{Proposal}

Draw a small group-valued perturbation \(R\in G\) and propose
\begin{equation}
  U_\mu'(x)=R\,U_\mu(x).
  \label{eq:proposal}
\end{equation}

\begin{interpretation}
Because both factors belong to \(G\), the proposal remains in \(G\). The
proposal distribution should be symmetric, or its forward/reverse probability
ratio must be included in a more general Metropolis--Hastings test.
\end{interpretation}

\subsection{Local action difference}

Only the \(2(D-1)\) plaquettes containing \(U_\mu(x)\) can change. If
\(\mathcal P(x,\mu)\) denotes that set, compute
\begin{align}
  \Delta S
  &=S_{\mathrm{local}}[U']-S_{\mathrm{local}}[U],
  \label{eq:delta-action}\\
  S_{\mathrm{local}}[U]
  &=\beta\sum_{p\in\mathcal P(x,\mu)}
  \left(1-\frac{1}{N}\ReTr U_p\right).
  \label{eq:local-action}
\end{align}

\begin{interpretation}
Equations~\eqref{eq:delta-action}--\eqref{eq:local-action} replace an
\(O(D^2\Vlat)\) full-action recomputation with a fixed-size local calculation.
As a correctness test, this \(\Delta S\) must agree with the difference of two
slow global-action evaluations.
\end{interpretation}

\subsection{Metropolis decision}

The acceptance probability is
\begin{equation}
  P_{\mathrm{accept}}=\min\left(1,\ee^{-\Delta S}\right).
  \label{eq:acceptance}
\end{equation}

For \(u\sim\operatorname{Uniform}[0,1)\), accept if
\begin{equation}
  u<P_{\mathrm{accept}},
  \qquad\text{equivalently}\qquad
  \log u<-\Delta S.
  \label{eq:log-acceptance}
\end{equation}

\begin{interpretation}
If \(\Delta S\leq0\), the proposal is always accepted. If \(\Delta S>0\),
it is accepted with probability \(\ee^{-\Delta S}\). The logarithmic form in
Equation~\eqref{eq:log-acceptance} avoids numerical underflow. Visiting every
site \(x\) and direction \(\mu\) once constitutes one sweep.
\end{interpretation}

\paragraph{Step output.}
One new Markov-chain state plus acceptance counters used to diagnose proposal
size and sampling behavior.

\section{Step 5: thermalize and measure}

\subsection{Sampling schedule}

First discard \(N_{\mathrm{warmup}}\) sweeps. Then save one measurement after
every \(N_{\mathrm{spacing}}\) sweeps until \(\Nmeas\) measurements have been
recorded. Saving less frequently can reduce storage and correlation, but it
does not substitute for estimating autocorrelation.

\subsection{Mean plaquette}

There are
\begin{equation}
  N_{\mathrm{plaq}}=\Vlat\frac{D(D-1)}{2}
  \label{eq:plaquette-count}
\end{equation}
unoriented plaquettes on a periodic isotropic lattice. Their normalized mean is
\begin{equation}
  \bar P=
  \frac{1}{N_{\mathrm{plaq}}}
  \sum_x\sum_{\mu<\nu}
  \frac{1}{N}\ReTr U_{\mu\nu}(x).
  \label{eq:mean-plaquette}
\end{equation}

\begin{interpretation}
\(\bar P\) is inexpensive and is the first physical sanity check. Its history
helps reveal thermalization and frozen or unstable updates, but it does not by
itself determine the static potential.
\end{interpretation}

\subsection{Wilson loops}

For a rectangular closed path \(C_{r,t}\), define
\begin{equation}
  W(r,t)=\frac{1}{N}
  \left\langle
    \ReTr\,
    \mathcal P\!\!\prod_{(x,\mu)\in C_{r,t}}U_\mu(x)
  \right\rangle.
  \label{eq:wilson-loop}
\end{equation}

\begin{interpretation}
\(\mathcal P\) means path ordering. Each forward segment contributes a link
\(U_\mu(x)\); each backward segment contributes its dagger. The trace closes
the color indices and makes the result gauge invariant. Measure several
spatial widths \(r\) and time extents \(t\) on every saved configuration.
\end{interpretation}

\paragraph{Step output.}
Correlated histories \(\bar P_k\) and \(W_k(r,t)\), each stored with chain ID
and sweep number.

\section{Step 6: extract the static potential}

A spectral decomposition has the schematic form
\begin{equation}
  W(r,t)=\sum_n |c_n(r)|^2\ee^{-aE_n(r)t}.
  \label{eq:spectral}
\end{equation}

\begin{interpretation}
The chosen loop operator overlaps with several energy eigenstates. Excited
states decay faster with \(t\), but the statistical signal also becomes noisier.
This creates the practical need for a fit window.
\end{interpretation}

At sufficiently large \(t\), the lowest state dominates:
\begin{equation}
  W(r,t)\simeq A(r)\ee^{-aV(r)t}.
  \label{eq:large-time}
\end{equation}

\begin{interpretation}
\(V(r)\) is the static-source potential and \(A(r)\) is an unknown overlap
factor. The product \(aV(r)\) is dimensionless when \(t\) is counted in lattice
time steps.
\end{interpretation}

The neighboring-time ratio removes \(A(r)\):
\begin{equation}
  aV_{\mathrm{eff}}(r,t)=
  \log\left(\frac{W(r,t)}{W(r,t+1)}\right).
  \label{eq:effective-potential}
\end{equation}

\begin{interpretation}
Plot Equation~\eqref{eq:effective-potential} against \(t\) at fixed \(r\).
A region consistent with a constant is a candidate ground-state plateau.
Fit that region to obtain \(aV(r)\), and vary the fit window to estimate a
systematic uncertainty.
\end{interpretation}

\paragraph{Step output.}
An effective-potential curve for each \(r\), a documented plateau window, and a
fitted \(aV(r)\) with uncertainty.

\section{Step 7: account for correlated samples}

For a measurement history \(O_1,\ldots,O_{\Nmeas}\), first compute
\begin{equation}
  \bar O=\frac{1}{\Nmeas}\sum_{k=1}^{\Nmeas}O_k.
  \label{eq:sample-mean}
\end{equation}

Define the lag-\(\tau\) autocovariance and normalized autocorrelation by
\begin{align}
  C_O(\tau)
  &=
  \frac{1}{\Nmeas-\tau}
  \sum_{k=1}^{\Nmeas-\tau}
  (O_k-\bar O)(O_{k+\tau}-\bar O),
  \label{eq:autocovariance}\\
  \rho_O(\tau)&=\frac{C_O(\tau)}{C_O(0)}.
  \label{eq:autocorrelation}
\end{align}

\begin{interpretation}
\(\rho_O(\tau)\) measures how strongly samples separated by \(\tau\) saved
measurements remain related. It equals one at zero lag and should decay toward
zero. Different observables can have different autocorrelation scales.
\end{interpretation}

With a justified truncation window \(\tau_{\max}\), estimate
\begin{equation}
  \tau_{\mathrm{int}}=
  \frac{1}{2}+
  \sum_{\tau=1}^{\tau_{\max}}\rho_O(\tau).
  \label{eq:integrated-autocorrelation}
\end{equation}

\begin{interpretation}
The sum cannot be extended indefinitely because large-lag estimates become
noise dominated. A windowing rule or a blocking study should justify
\(\tau_{\max}\).
\end{interpretation}

The approximate effective sample count and variance of the mean are
\begin{align}
  N_{\mathrm{eff}}
  &\simeq\frac{\Nmeas}{2\tau_{\mathrm{int}}},
  \label{eq:effective-samples}\\
  \operatorname{Var}(\bar O)
  &\simeq
  \frac{2\tau_{\mathrm{int}}}{\Nmeas}
  \operatorname{Var}(O).
  \label{eq:correlated-variance}
\end{align}

\begin{interpretation}
If the samples were independent, \(\tau_{\mathrm{int}}\approx1/2\) and
\(N_{\mathrm{eff}}\approx\Nmeas\). Positive correlation reduces the effective
sample count and enlarges the error bar. Wilson-loop ratios and potential fits
must be recomputed inside each block-jackknife or bootstrap resample so their
cross-correlations are preserved.
\end{interpretation}

\paragraph{Step output.}
Autocorrelation estimates, a stable block size, covariance-aware uncertainties,
and an effective sample count for each observable.

\section{Step 8: validation gates}

The following checks should be passed in order:

\begin{enumerate}
  \item coordinate-to-index round trips and every periodic boundary;
  \item unitarity and determinant constraints for generated links;
  \item local \(\Delta S\) versus a recomputed global action difference;
  \item invariance of traced closed loops under random gauge transformations;
  \item agreement of hot and cold starts after thermalization;
  \item agreement of independent chains within correlated uncertainties;
  \item a tiny-lattice calculation or published benchmark with matching
        conventions; and
  \item volume and lattice-spacing studies before making physical claims.
\end{enumerate}

A failed gate does not necessarily make the raw run useless for debugging, but
it does make the resulting observable scientifically unvalidated.

\paragraph{Step output.}
A validation record stating which structural, numerical, statistical, and
physical checks passed or failed.

\section{Reference algorithm}

\begin{verbatim}
validate configuration
construct lattice indices and periodic neighbors
initialize all links and deterministic random streams

repeat N_warmup times:
    sweep over every site x and direction mu

repeat N_measurements times:
    repeat N_spacing times:
        sweep over every site x and direction mu
    measure mean plaquette and Wilson loops W(r,t)

block the correlated measurements
for every jackknife or bootstrap resample:
    recompute W(r,t)
    compute a V_eff(r,t) = log[W(r,t) / W(r,t+1)]
    fit the selected plateau

report means, covariance-aware uncertainty, diagnostics, and provenance
\end{verbatim}

\section{How to interpret the final gauge result}

A reported value of \(aV(r)\) is not automatically a continuum QCD prediction.
It is a result for the documented action, coupling, volume, boundary conditions,
operator, sampling procedure, and fit window. A defensible report therefore
includes the resolved configuration, action convention, chain histories,
autocorrelation treatment, validation results, and known finite-volume and
lattice-spacing limitations.

\end{document}
